% Options for packages loaded elsewhere
\PassOptionsToPackage{unicode}{hyperref}
\PassOptionsToPackage{hyphens}{url}
\documentclass[
]{article}
\usepackage{xcolor}
\usepackage[margin=1in]{geometry}
\usepackage{amsmath,amssymb}
\setcounter{secnumdepth}{-\maxdimen} % remove section numbering
\usepackage{iftex}
\ifPDFTeX
  \usepackage[T1]{fontenc}
  \usepackage[utf8]{inputenc}
  \usepackage{textcomp} % provide euro and other symbols
\else % if luatex or xetex
  \usepackage{unicode-math} % this also loads fontspec
  \defaultfontfeatures{Scale=MatchLowercase}
  \defaultfontfeatures[\rmfamily]{Ligatures=TeX,Scale=1}
\fi
\usepackage{lmodern}
\ifPDFTeX\else
  % xetex/luatex font selection
\fi
% Use upquote if available, for straight quotes in verbatim environments
\IfFileExists{upquote.sty}{\usepackage{upquote}}{}
\IfFileExists{microtype.sty}{% use microtype if available
  \usepackage[]{microtype}
  \UseMicrotypeSet[protrusion]{basicmath} % disable protrusion for tt fonts
}{}
\makeatletter
\@ifundefined{KOMAClassName}{% if non-KOMA class
  \IfFileExists{parskip.sty}{%
    \usepackage{parskip}
  }{% else
    \setlength{\parindent}{0pt}
    \setlength{\parskip}{6pt plus 2pt minus 1pt}}
}{% if KOMA class
  \KOMAoptions{parskip=half}}
\makeatother
\usepackage{color}
\usepackage{fancyvrb}
\newcommand{\VerbBar}{|}
\newcommand{\VERB}{\Verb[commandchars=\\\{\}]}
\DefineVerbatimEnvironment{Highlighting}{Verbatim}{commandchars=\\\{\}}
% Add ',fontsize=\small' for more characters per line
\usepackage{framed}
\definecolor{shadecolor}{RGB}{248,248,248}
\newenvironment{Shaded}{\begin{snugshade}}{\end{snugshade}}
\newcommand{\AlertTok}[1]{\textcolor[rgb]{0.94,0.16,0.16}{#1}}
\newcommand{\AnnotationTok}[1]{\textcolor[rgb]{0.56,0.35,0.01}{\textbf{\textit{#1}}}}
\newcommand{\AttributeTok}[1]{\textcolor[rgb]{0.13,0.29,0.53}{#1}}
\newcommand{\BaseNTok}[1]{\textcolor[rgb]{0.00,0.00,0.81}{#1}}
\newcommand{\BuiltInTok}[1]{#1}
\newcommand{\CharTok}[1]{\textcolor[rgb]{0.31,0.60,0.02}{#1}}
\newcommand{\CommentTok}[1]{\textcolor[rgb]{0.56,0.35,0.01}{\textit{#1}}}
\newcommand{\CommentVarTok}[1]{\textcolor[rgb]{0.56,0.35,0.01}{\textbf{\textit{#1}}}}
\newcommand{\ConstantTok}[1]{\textcolor[rgb]{0.56,0.35,0.01}{#1}}
\newcommand{\ControlFlowTok}[1]{\textcolor[rgb]{0.13,0.29,0.53}{\textbf{#1}}}
\newcommand{\DataTypeTok}[1]{\textcolor[rgb]{0.13,0.29,0.53}{#1}}
\newcommand{\DecValTok}[1]{\textcolor[rgb]{0.00,0.00,0.81}{#1}}
\newcommand{\DocumentationTok}[1]{\textcolor[rgb]{0.56,0.35,0.01}{\textbf{\textit{#1}}}}
\newcommand{\ErrorTok}[1]{\textcolor[rgb]{0.64,0.00,0.00}{\textbf{#1}}}
\newcommand{\ExtensionTok}[1]{#1}
\newcommand{\FloatTok}[1]{\textcolor[rgb]{0.00,0.00,0.81}{#1}}
\newcommand{\FunctionTok}[1]{\textcolor[rgb]{0.13,0.29,0.53}{\textbf{#1}}}
\newcommand{\ImportTok}[1]{#1}
\newcommand{\InformationTok}[1]{\textcolor[rgb]{0.56,0.35,0.01}{\textbf{\textit{#1}}}}
\newcommand{\KeywordTok}[1]{\textcolor[rgb]{0.13,0.29,0.53}{\textbf{#1}}}
\newcommand{\NormalTok}[1]{#1}
\newcommand{\OperatorTok}[1]{\textcolor[rgb]{0.81,0.36,0.00}{\textbf{#1}}}
\newcommand{\OtherTok}[1]{\textcolor[rgb]{0.56,0.35,0.01}{#1}}
\newcommand{\PreprocessorTok}[1]{\textcolor[rgb]{0.56,0.35,0.01}{\textit{#1}}}
\newcommand{\RegionMarkerTok}[1]{#1}
\newcommand{\SpecialCharTok}[1]{\textcolor[rgb]{0.81,0.36,0.00}{\textbf{#1}}}
\newcommand{\SpecialStringTok}[1]{\textcolor[rgb]{0.31,0.60,0.02}{#1}}
\newcommand{\StringTok}[1]{\textcolor[rgb]{0.31,0.60,0.02}{#1}}
\newcommand{\VariableTok}[1]{\textcolor[rgb]{0.00,0.00,0.00}{#1}}
\newcommand{\VerbatimStringTok}[1]{\textcolor[rgb]{0.31,0.60,0.02}{#1}}
\newcommand{\WarningTok}[1]{\textcolor[rgb]{0.56,0.35,0.01}{\textbf{\textit{#1}}}}
\usepackage{graphicx}
\makeatletter
\newsavebox\pandoc@box
\newcommand*\pandocbounded[1]{% scales image to fit in text height/width
  \sbox\pandoc@box{#1}%
  \Gscale@div\@tempa{\textheight}{\dimexpr\ht\pandoc@box+\dp\pandoc@box\relax}%
  \Gscale@div\@tempb{\linewidth}{\wd\pandoc@box}%
  \ifdim\@tempb\p@<\@tempa\p@\let\@tempa\@tempb\fi% select the smaller of both
  \ifdim\@tempa\p@<\p@\scalebox{\@tempa}{\usebox\pandoc@box}%
  \else\usebox{\pandoc@box}%
  \fi%
}
% Set default figure placement to htbp
\def\fps@figure{htbp}
\makeatother
\setlength{\emergencystretch}{3em} % prevent overfull lines
\providecommand{\tightlist}{%
  \setlength{\itemsep}{0pt}\setlength{\parskip}{0pt}}
\usepackage{bookmark}
\IfFileExists{xurl.sty}{\usepackage{xurl}}{} % add URL line breaks if available
\urlstyle{same}
% fallback for those not using the hyperref driver hyperxmp:
\makeatletter
\@ifundefined{xmpquote}{\newcommand{\xmpquote}[1]{#1}}{}
\makeatother
\hypersetup{
  pdftitle={Assignment 1},
  hidelinks,
  pdfcreator={LaTeX via pandoc}}

\title{Assignment 1}
\author{}
\date{\vspace{-2.5em}}

\begin{document}
\maketitle

\section{Task 1: Literature}\label{task-1-literature}

\begin{quote}
Developmental landscape of human forebrain at a single-cell level
identifies early waves of oligodendrogenesis. David van Bruggen, et
al.~PMID: 35523173
\end{quote}

\subsubsection{Q1. What is the medically relevant insight from the
article?}\label{q1.-what-is-the-medically-relevant-insight-from-the-article}

This study doesn't focus on any medical trait, and thus the medical
relevance is not mentioned in the paper. However, oligodendrocyte
precursor cells have been linked to psychiatric disorders like
schizophrenia and bipolar disorder, so findings from this study might
help understand the neural mechanisms of mental disorders.

\subsubsection{Q2. Which genomics technology/ technologies were
used?}\label{q2.-which-genomics-technology-technologies-were-used}

single-cell RNA sequencing: examining the identity of cells based on
gene expression profiles

single-cell ATAC-seq: examining the chromatin accessibility, which is
required for oligodentrogenesis

Spatial \emph{in situ} sequencing: locating the sites for early human
oligodendrogenesis

RNAscope HiPlex \emph{in situ} hybridization: validating findings
Spatial ISS analysis

\subsubsection{Q3. List and explain at least three questions/ hypotheses
you can think of that extend the analysis presented in the
paper.}\label{q3.-list-and-explain-at-least-three-questions-hypotheses-you-can-think-of-that-extend-the-analysis-presented-in-the-paper.}

\begin{enumerate}
\def\labelenumi{\arabic{enumi}.}
\tightlist
\item
  How does oligodendrogenesis progress in midbrain and hidbrain
  (compared to forebrain)?
\item
  If the first wave of oligodendrogenesis is evolutionary conservative,
  are any later oligodendrogenesis waves less evolutionary conservative
  (i.e., more specific to humans)?
\item
  If the 2nd is true, are those less evolutionary conservative
  oligodendrogenesis waves with more individual differences (in for
  example time), and related to human-specific cognitive traits like
  intelligence and psychiatric disorders?
\end{enumerate}

\subsubsection{Q4. {[}Optional{]} Devise a computational analysis
strategy for (some of) the listed questions under
3a.}\label{q4.-optional-devise-a-computational-analysis-strategy-for-some-of-the-listed-questions-under-3a.}

\section{Task 4: R basic operations}\label{task-4-r-basic-operations}

\begin{Shaded}
\begin{Highlighting}[]
\FunctionTok{cat}\NormalTok{(}\StringTok{\textquotesingle{}1. Square root of 10:\textquotesingle{}}\NormalTok{, }\FunctionTok{sqrt}\NormalTok{(}\DecValTok{10}\NormalTok{), }\StringTok{\textquotesingle{}}\SpecialCharTok{\textbackslash{}n}\StringTok{\textquotesingle{}}\NormalTok{)}
\end{Highlighting}
\end{Shaded}

\begin{verbatim}
## 1. Square root of 10: 3.162278
\end{verbatim}

\begin{Shaded}
\begin{Highlighting}[]
\FunctionTok{cat}\NormalTok{(}\StringTok{\textquotesingle{}2. Logarithm of 32 to the base 2:\textquotesingle{}}\NormalTok{, }\FunctionTok{log2}\NormalTok{(}\DecValTok{32}\NormalTok{), }\StringTok{\textquotesingle{}}\SpecialCharTok{\textbackslash{}n}\StringTok{\textquotesingle{}}\NormalTok{)}
\end{Highlighting}
\end{Shaded}

\begin{verbatim}
## 2. Logarithm of 32 to the base 2: 5
\end{verbatim}

\begin{Shaded}
\begin{Highlighting}[]
\FunctionTok{cat}\NormalTok{(}\StringTok{\textquotesingle{}3. the sum of the numbers from 1 to 1000:\textquotesingle{}}\NormalTok{, }\FunctionTok{sum}\NormalTok{(}\FunctionTok{seq}\NormalTok{(}\DecValTok{1000}\NormalTok{)), }\StringTok{\textquotesingle{}}\SpecialCharTok{\textbackslash{}n}\StringTok{\textquotesingle{}}\NormalTok{)}
\end{Highlighting}
\end{Shaded}

\begin{verbatim}
## 3. the sum of the numbers from 1 to 1000: 500500
\end{verbatim}

\begin{Shaded}
\begin{Highlighting}[]
\FunctionTok{cat}\NormalTok{(}\StringTok{\textquotesingle{}4. the sum of all even numbers from 2 to 1000:\textquotesingle{}}\NormalTok{, }\FunctionTok{sum}\NormalTok{(}\FunctionTok{seq}\NormalTok{(}\DecValTok{2}\NormalTok{, }\DecValTok{1000}\NormalTok{, }\DecValTok{2}\NormalTok{)), }\StringTok{\textquotesingle{}}\SpecialCharTok{\textbackslash{}n}\StringTok{\textquotesingle{}}\NormalTok{)}
\end{Highlighting}
\end{Shaded}

\begin{verbatim}
## 4. the sum of all even numbers from 2 to 1000: 250500
\end{verbatim}

\begin{Shaded}
\begin{Highlighting}[]
\FunctionTok{cat}\NormalTok{(}\StringTok{\textquotesingle{}5. the number of pairwise comparisons for 100 genes:\textquotesingle{}}\NormalTok{, (}\DecValTok{100}\SpecialCharTok{*}\DecValTok{100{-}100}\NormalTok{)}\SpecialCharTok{/}\DecValTok{2}\NormalTok{, }\StringTok{\textquotesingle{}}\SpecialCharTok{\textbackslash{}n}\StringTok{\textquotesingle{}}\NormalTok{)}
\end{Highlighting}
\end{Shaded}

\begin{verbatim}
## 5. the number of pairwise comparisons for 100 genes: 4950
\end{verbatim}

\begin{Shaded}
\begin{Highlighting}[]
\FunctionTok{cat}\NormalTok{(}\StringTok{\textquotesingle{}6. the number of ways to arrange 100 genes in triplets:\textquotesingle{}}\NormalTok{, }\FunctionTok{factorial}\NormalTok{(}\DecValTok{100}\NormalTok{)}\SpecialCharTok{/}\FunctionTok{factorial}\NormalTok{(}\DecValTok{100{-}3}\NormalTok{), }\StringTok{\textquotesingle{}}\SpecialCharTok{\textbackslash{}n}\StringTok{\textquotesingle{}}\NormalTok{)}
\end{Highlighting}
\end{Shaded}

\begin{verbatim}
## 6. the number of ways to arrange 100 genes in triplets: 970200
\end{verbatim}

\section{Task 5: Using R example
datasets}\label{task-5-using-r-example-datasets}

\subsubsection{1. Load the CO2 dataset}\label{load-the-co2-dataset}

\begin{Shaded}
\begin{Highlighting}[]
\FunctionTok{data}\NormalTok{(}\StringTok{\textquotesingle{}CO2\textquotesingle{}}\NormalTok{)}
\end{Highlighting}
\end{Shaded}

\subsubsection{2. Describe the dataset}\label{describe-the-dataset}

\begin{Shaded}
\begin{Highlighting}[]
\FunctionTok{help}\NormalTok{(CO2)}
\FunctionTok{cat}\NormalTok{(}\StringTok{\textquotesingle{}Description:}\SpecialCharTok{\textbackslash{}n}\StringTok{\textquotesingle{}}\NormalTok{, }\StringTok{\textquotesingle{}The CO2 data frame has 84 rows and 5 columns of data from an experiment on the cold tolerance of the grass species Echinochloa crus{-}galli.\textquotesingle{}}\NormalTok{)}
\end{Highlighting}
\end{Shaded}

\begin{verbatim}
## Description:
##  The CO2 data frame has 84 rows and 5 columns of data from an experiment on the cold tolerance of the grass species Echinochloa crus-galli.
\end{verbatim}

\subsubsection{3. the average and median CO2 uptake of the plants from
Quebec and
Mississippi}\label{the-average-and-median-co2-uptake-of-the-plants-from-quebec-and-mississippi}

\begin{Shaded}
\begin{Highlighting}[]
\NormalTok{summary\_uptake }\OtherTok{=} \FunctionTok{data.frame}\NormalTok{()}
\NormalTok{cols\_summary }\OtherTok{=} \FunctionTok{c}\NormalTok{(}\StringTok{\textquotesingle{}type\textquotesingle{}}\NormalTok{, }\StringTok{\textquotesingle{}mean\textquotesingle{}}\NormalTok{, }\StringTok{\textquotesingle{}median\textquotesingle{}}\NormalTok{)}
\NormalTok{types }\OtherTok{=} \FunctionTok{c}\NormalTok{(}\StringTok{\textquotesingle{}Quebec\textquotesingle{}}\NormalTok{, }\StringTok{\textquotesingle{}Mississippi\textquotesingle{}}\NormalTok{)}

\CommentTok{\# summary of each type}
\ControlFlowTok{for}\NormalTok{ (i }\ControlFlowTok{in} \DecValTok{1}\SpecialCharTok{:}\FunctionTok{length}\NormalTok{(types)) \{}
\NormalTok{  type }\OtherTok{=}\NormalTok{ types[i]}
\NormalTok{  uptake\_type }\OtherTok{=} \FunctionTok{na.exclude}\NormalTok{(CO2[CO2}\SpecialCharTok{$}\NormalTok{Type }\SpecialCharTok{==}\NormalTok{ type, }\StringTok{\textquotesingle{}uptake\textquotesingle{}}\NormalTok{])}
  
\NormalTok{  summary\_uptake[i, cols\_summary] }\OtherTok{=} \FunctionTok{c}\NormalTok{(type,}
                                      \FunctionTok{round}\NormalTok{(}\FunctionTok{mean}\NormalTok{(uptake\_type), }\DecValTok{2}\NormalTok{),}
                                      \FunctionTok{round}\NormalTok{(}\FunctionTok{median}\NormalTok{(uptake\_type), }\DecValTok{2}\NormalTok{))\}}
  
\FunctionTok{print}\NormalTok{(summary\_uptake)}
\end{Highlighting}
\end{Shaded}

\begin{verbatim}
##          type  mean median
## 1      Quebec 33.54  37.15
## 2 Mississippi 20.88   19.3
\end{verbatim}

\subsubsection{4. {[}optional{]} genes expressed in each sample; genes
not expressed in any
sample}\label{optional-genes-expressed-in-each-sample-genes-not-expressed-in-any-sample}

\begin{Shaded}
\begin{Highlighting}[]
\CommentTok{\# BiocManager::install("airway")}
\CommentTok{\# library(\textquotesingle{}airway\textquotesingle{})}
\CommentTok{\# data(\textquotesingle{}airway\textquotesingle{})}
\end{Highlighting}
\end{Shaded}

\section{Task 6: R Functions}\label{task-6-r-functions}

\subsubsection{1. A function calculating the ratio of the mean and the
median of a
vector}\label{a-function-calculating-the-ratio-of-the-mean-and-the-median-of-a-vector}

\begin{Shaded}
\begin{Highlighting}[]
\NormalTok{mean\_over\_median }\OtherTok{=} \ControlFlowTok{function}\NormalTok{(x) \{}
\NormalTok{  x }\OtherTok{=} \FunctionTok{na.exclude}\NormalTok{(x)}
\NormalTok{  ratio }\OtherTok{=} \FunctionTok{mean}\NormalTok{(x)}\SpecialCharTok{/}\FunctionTok{median}\NormalTok{(x)}
  \FunctionTok{return}\NormalTok{(ratio)}
\NormalTok{\}}

\CommentTok{\# test}
\FunctionTok{print}\NormalTok{(}\FunctionTok{mean\_over\_median}\NormalTok{(}\FunctionTok{c}\NormalTok{(}\DecValTok{2}\NormalTok{, }\DecValTok{2}\NormalTok{, }\DecValTok{3}\NormalTok{, }\DecValTok{4}\NormalTok{, }\DecValTok{5}\NormalTok{)))}
\end{Highlighting}
\end{Shaded}

\begin{verbatim}
## [1] 1.066667
\end{verbatim}

\subsubsection{2. A function calculating the mean after discarding the
min and the
max}\label{a-function-calculating-the-mean-after-discarding-the-min-and-the-max}

\begin{Shaded}
\begin{Highlighting}[]
\NormalTok{mean\_without\_extreme }\OtherTok{=} \ControlFlowTok{function}\NormalTok{(x) \{}
\NormalTok{  x }\OtherTok{=} \FunctionTok{na.exclude}\NormalTok{(x)}
\NormalTok{  x }\OtherTok{=}\NormalTok{ x[(x }\SpecialCharTok{!=} \FunctionTok{min}\NormalTok{(x)) }\SpecialCharTok{\&}\NormalTok{ (x }\SpecialCharTok{!=} \FunctionTok{max}\NormalTok{(x))]}
  
  \FunctionTok{return}\NormalTok{(}\FunctionTok{mean}\NormalTok{(x))}
\NormalTok{\}}

\CommentTok{\# test}
\FunctionTok{mean\_without\_extreme}\NormalTok{(}\FunctionTok{c}\NormalTok{(}\DecValTok{1}\NormalTok{, }\DecValTok{2}\NormalTok{, }\DecValTok{2}\NormalTok{, }\DecValTok{5}\NormalTok{))}
\end{Highlighting}
\end{Shaded}

\begin{verbatim}
## [1] 2
\end{verbatim}

\subsubsection{3. Explaining ``piping''}\label{explaining-piping}

A pipe is discouraged when the included functions (1) have multiple
inputs or outputs, (2) have a complex dependency, (3) use the current
environment, or (4) use lazy evaluation. Besides, (5) piping too many
intermediate steps/functions can make the code harder to read and debug.

\subsubsection{4. Explaining ``apply-family
functions''}\label{explaining-apply-family-functions}

The apply-family functions are used to apply a function to each element
in an iterable object. Though they won't make the running more
efficient, they replace a for-loop with a more compact code snippet. In
other words, it seems like they only make the code look better.

\section{Task 7: Basic visualization with
R}\label{task-7-basic-visualization-with-r}

\subsubsection{1. Compare the distributions of the body heights of the
two species from the `magic\_guys.csv' dataset
graphically}\label{compare-the-distributions-of-the-body-heights-of-the-two-species-from-the-magic_guys.csv-dataset-graphically}

\begin{Shaded}
\begin{Highlighting}[]
\FunctionTok{library}\NormalTok{(ggplot2)}
\CommentTok{\# with basic hist function: Length of Jedi}
\NormalTok{magic }\OtherTok{=} \FunctionTok{read.csv}\NormalTok{(}\StringTok{\textquotesingle{}data/magic\_guys.csv\textquotesingle{}}\NormalTok{, }\AttributeTok{header=}\NormalTok{T)}
\FunctionTok{hist}\NormalTok{(magic[magic}\SpecialCharTok{$}\NormalTok{species }\SpecialCharTok{==} \StringTok{\textquotesingle{}jedi\textquotesingle{}}\NormalTok{, }\StringTok{\textquotesingle{}length\textquotesingle{}}\NormalTok{], }\AttributeTok{breaks=}\DecValTok{8}\NormalTok{, }\AttributeTok{xlim=}\FunctionTok{c}\NormalTok{(}\DecValTok{0}\NormalTok{, }\DecValTok{300}\NormalTok{),}
     \AttributeTok{xlab=}\StringTok{\textquotesingle{}Length of Jedi\textquotesingle{}}\NormalTok{, }\AttributeTok{main=}\StringTok{\textquotesingle{}\textquotesingle{}}\NormalTok{)}
\end{Highlighting}
\end{Shaded}

\pandocbounded{\includegraphics[keepaspectratio]{assignment_1_files/figure-latex/unnamed-chunk-8-1.pdf}}

\begin{Shaded}
\begin{Highlighting}[]
\CommentTok{\# with basic hist function: Length of Sith}
\FunctionTok{hist}\NormalTok{(magic[magic}\SpecialCharTok{$}\NormalTok{species }\SpecialCharTok{==} \StringTok{\textquotesingle{}sith\textquotesingle{}}\NormalTok{, }\StringTok{\textquotesingle{}length\textquotesingle{}}\NormalTok{], }\AttributeTok{breaks=}\DecValTok{8}\NormalTok{, }\AttributeTok{xlim=}\FunctionTok{c}\NormalTok{(}\DecValTok{0}\NormalTok{, }\DecValTok{300}\NormalTok{),}
     \AttributeTok{xlab=}\StringTok{\textquotesingle{}Length of Sith\textquotesingle{}}\NormalTok{, }\AttributeTok{main=}\StringTok{\textquotesingle{}\textquotesingle{}}\NormalTok{)}
\end{Highlighting}
\end{Shaded}

\pandocbounded{\includegraphics[keepaspectratio]{assignment_1_files/figure-latex/unnamed-chunk-8-2.pdf}}

\begin{Shaded}
\begin{Highlighting}[]
\CommentTok{\# with ggplot and geom\_histogram}
\NormalTok{n\_bins }\OtherTok{=} \DecValTok{10}
\NormalTok{hist\_length }\OtherTok{=} \FunctionTok{ggplot}\NormalTok{(magic, }\FunctionTok{aes}\NormalTok{(}\AttributeTok{x=}\NormalTok{length, }\AttributeTok{fill=}\NormalTok{species)) }\SpecialCharTok{+}
  \FunctionTok{geom\_histogram}\NormalTok{(}\AttributeTok{bins=}\NormalTok{n\_bins, }\AttributeTok{alpha=}\FloatTok{0.5}\NormalTok{, }\AttributeTok{position=}\StringTok{\textquotesingle{}identity\textquotesingle{}}\NormalTok{)}

\FunctionTok{ggsave}\NormalTok{(}\StringTok{"results/histogram\_height.pdf"}\NormalTok{, }\AttributeTok{plot=}\NormalTok{hist\_length)}
\end{Highlighting}
\end{Shaded}

\begin{verbatim}
## Saving 6.5 x 4.5 in image
\end{verbatim}

\begin{Shaded}
\begin{Highlighting}[]
\FunctionTok{ggsave}\NormalTok{(}\StringTok{"results/histogram\_height.png"}\NormalTok{, }\AttributeTok{plot=}\NormalTok{hist\_length)}
\end{Highlighting}
\end{Shaded}

\begin{verbatim}
## Saving 6.5 x 4.5 in image
\end{verbatim}

\begin{Shaded}
\begin{Highlighting}[]
\FunctionTok{ggsave}\NormalTok{(}\StringTok{"results/histogram\_height.svg"}\NormalTok{, }\AttributeTok{plot=}\NormalTok{hist\_length)}
\end{Highlighting}
\end{Shaded}

\begin{verbatim}
## Saving 6.5 x 4.5 in image
\end{verbatim}

\begin{Shaded}
\begin{Highlighting}[]
\FunctionTok{print}\NormalTok{(hist\_length)}
\end{Highlighting}
\end{Shaded}

\pandocbounded{\includegraphics[keepaspectratio]{assignment_1_files/figure-latex/unnamed-chunk-8-3.pdf}}

\begin{Shaded}
\begin{Highlighting}[]
\NormalTok{box\_height }\OtherTok{=} \FunctionTok{ggplot}\NormalTok{(magic, }\FunctionTok{aes}\NormalTok{(}\AttributeTok{x=}\NormalTok{species, }\AttributeTok{y=}\NormalTok{length)) }\SpecialCharTok{+}
  \FunctionTok{geom\_boxplot}\NormalTok{()}
\FunctionTok{ggsave}\NormalTok{(}\StringTok{"results/box\_height.pdf"}\NormalTok{, }\AttributeTok{plot=}\NormalTok{box\_height)}
\end{Highlighting}
\end{Shaded}

\begin{verbatim}
## Saving 6.5 x 4.5 in image
\end{verbatim}

\begin{Shaded}
\begin{Highlighting}[]
\FunctionTok{ggsave}\NormalTok{(}\StringTok{"results/box\_height.png"}\NormalTok{, }\AttributeTok{plot=}\NormalTok{box\_height)}
\end{Highlighting}
\end{Shaded}

\begin{verbatim}
## Saving 6.5 x 4.5 in image
\end{verbatim}

\begin{Shaded}
\begin{Highlighting}[]
\FunctionTok{ggsave}\NormalTok{(}\StringTok{"results/box\_height.svg"}\NormalTok{, }\AttributeTok{plot=}\NormalTok{box\_height)}
\end{Highlighting}
\end{Shaded}

\begin{verbatim}
## Saving 6.5 x 4.5 in image
\end{verbatim}

\begin{Shaded}
\begin{Highlighting}[]
\FunctionTok{print}\NormalTok{(box\_height)}
\end{Highlighting}
\end{Shaded}

\pandocbounded{\includegraphics[keepaspectratio]{assignment_1_files/figure-latex/unnamed-chunk-9-1.pdf}}

Formats of `pdf' and `svg' are usually preferred, as they are vector
graphics, which has high image quality, small file size, while allows
for editing later.

\subsubsection{2. Inspect and visualize
`microarray\_data.tab'}\label{inspect-and-visualize-microarray_data.tab}

\begin{Shaded}
\begin{Highlighting}[]
\NormalTok{microarray }\OtherTok{=} \FunctionTok{read.table}\NormalTok{(}\StringTok{\textquotesingle{}data/microarray\_data.tab\textquotesingle{}}\NormalTok{, }\AttributeTok{sep=}\StringTok{\textquotesingle{}}\SpecialCharTok{\textbackslash{}t}\StringTok{\textquotesingle{}}\NormalTok{, }\AttributeTok{header=}\NormalTok{T)}
\CommentTok{\# a. inspect the size of the data frame}
\FunctionTok{print}\NormalTok{(}\FunctionTok{sprintf}\NormalTok{(}\StringTok{\textquotesingle{}The matrix has \%s rows and \%s columns\textquotesingle{}}\NormalTok{, }\FunctionTok{nrow}\NormalTok{(microarray), }\FunctionTok{ncol}\NormalTok{(microarray)))}
\end{Highlighting}
\end{Shaded}

\begin{verbatim}
## [1] "The matrix has 553 rows and 1000 columns"
\end{verbatim}

\begin{Shaded}
\begin{Highlighting}[]
\CommentTok{\# b. visualize the number of missing values per gene}
\NormalTok{counts\_missing }\OtherTok{=} \FunctionTok{apply}\NormalTok{(}\FunctionTok{is.na}\NormalTok{(microarray), }\DecValTok{2}\NormalTok{, sum)}
\FunctionTok{names}\NormalTok{(counts\_missing) }\OtherTok{=} \FunctionTok{as.numeric}\NormalTok{(}\FunctionTok{gsub}\NormalTok{(}\StringTok{\textquotesingle{}g\textquotesingle{}}\NormalTok{, }\StringTok{\textquotesingle{}\textquotesingle{}}\NormalTok{, }\FunctionTok{names}\NormalTok{(counts\_missing)))}
\FunctionTok{barplot}\NormalTok{(counts\_missing, }\AttributeTok{xlab=}\StringTok{\textquotesingle{}Gene Number\textquotesingle{}}\NormalTok{)}
\end{Highlighting}
\end{Shaded}

\pandocbounded{\includegraphics[keepaspectratio]{assignment_1_files/figure-latex/unnamed-chunk-11-1.pdf}}

\begin{Shaded}
\begin{Highlighting}[]
\CommentTok{\# c. find genes with specified missing rates}
\CommentTok{\# define the function for searching}
\NormalTok{find\_missing\_genes }\OtherTok{=} \ControlFlowTok{function}\NormalTok{(df, miss\_rate) \{}
\NormalTok{  counts\_missing  }\OtherTok{=} \FunctionTok{apply}\NormalTok{(}\FunctionTok{is.na}\NormalTok{(df), }\DecValTok{2}\NormalTok{, sum)}
\NormalTok{  percent\_missing }\OtherTok{=}\NormalTok{ counts\_missing}\SpecialCharTok{/}\FunctionTok{nrow}\NormalTok{(df)}
  
\NormalTok{  list\_result\_genes }\OtherTok{=} \FunctionTok{names}\NormalTok{(percent\_missing[percent\_missing }\SpecialCharTok{\textgreater{}}\NormalTok{ miss\_rate])}
  \FunctionTok{return}\NormalTok{(list\_result\_genes)}
\NormalTok{\}}

\CommentTok{\# apply the function and report the results}
\ControlFlowTok{for}\NormalTok{ (rate }\ControlFlowTok{in} \FunctionTok{c}\NormalTok{(}\FloatTok{0.1}\NormalTok{, }\FloatTok{0.2}\NormalTok{, }\FloatTok{0.5}\NormalTok{)) \{}
\NormalTok{  num\_genes }\OtherTok{=} \FunctionTok{length}\NormalTok{(}\FunctionTok{find\_missing\_genes}\NormalTok{(microarray, rate))}
  \FunctionTok{print}\NormalTok{(}\FunctionTok{sprintf}\NormalTok{(}\StringTok{\textquotesingle{}\%s genes with missing rate \textgreater{} \%s\textquotesingle{}}\NormalTok{, num\_genes, rate))}
\NormalTok{\}}
\end{Highlighting}
\end{Shaded}

\begin{verbatim}
## [1] "529 genes with missing rate > 0.1"
## [1] "377 genes with missing rate > 0.2"
## [1] "74 genes with missing rate > 0.5"
\end{verbatim}

\begin{Shaded}
\begin{Highlighting}[]
\CommentTok{\# d. replacing missing values with the average expression}
\ControlFlowTok{for}\NormalTok{ (col }\ControlFlowTok{in} \FunctionTok{names}\NormalTok{(microarray)) \{}
\NormalTok{  microarray[}\FunctionTok{is.na}\NormalTok{(microarray[, col]), col] }\OtherTok{=} \FunctionTok{mean}\NormalTok{(microarray[, col], }\AttributeTok{na.rm=}\NormalTok{T)}
\NormalTok{\}}
\end{Highlighting}
\end{Shaded}

\subsubsection{3. Visualize the CO2 data}\label{visualize-the-co2-data}

\begin{Shaded}
\begin{Highlighting}[]
\FunctionTok{ggplot}\NormalTok{(CO2, }\FunctionTok{aes}\NormalTok{(}\AttributeTok{x=}\NormalTok{Type, }\AttributeTok{y=}\NormalTok{uptake, }\AttributeTok{fill=}\NormalTok{Treatment)) }\SpecialCharTok{+}
  \FunctionTok{geom\_boxplot}\NormalTok{()}
\end{Highlighting}
\end{Shaded}

\pandocbounded{\includegraphics[keepaspectratio]{assignment_1_files/figure-latex/unnamed-chunk-14-1.pdf}}

\begin{Shaded}
\begin{Highlighting}[]
\CommentTok{\# merge define subgroups based on Type and Treatment}
\NormalTok{CO2[, }\StringTok{\textquotesingle{}Group\textquotesingle{}}\NormalTok{] }\OtherTok{=} \FunctionTok{paste}\NormalTok{(CO2[, }\StringTok{\textquotesingle{}Type\textquotesingle{}}\NormalTok{], CO2[, }\StringTok{\textquotesingle{}Treatment\textquotesingle{}}\NormalTok{], }\AttributeTok{sep=}\StringTok{\textquotesingle{}{-}\textquotesingle{}}\NormalTok{)}
\FunctionTok{ggplot}\NormalTok{(CO2, }\FunctionTok{aes}\NormalTok{(}\AttributeTok{x=}\NormalTok{conc, }\AttributeTok{y=}\NormalTok{uptake)) }\SpecialCharTok{+}
  \FunctionTok{geom\_point}\NormalTok{(}\FunctionTok{aes}\NormalTok{(}\AttributeTok{color=}\NormalTok{Group), }\AttributeTok{size =} \DecValTok{2}\NormalTok{) }\SpecialCharTok{+}
  \FunctionTok{facet\_grid}\NormalTok{(Type }\SpecialCharTok{\textasciitilde{}}\NormalTok{ Treatment)}
\end{Highlighting}
\end{Shaded}

\pandocbounded{\includegraphics[keepaspectratio]{assignment_1_files/figure-latex/unnamed-chunk-14-2.pdf}}

The bar plot shows that the Quebec types uptakes more CO2 than the
Mississippi type, and the chilled treatment decreases CO2 uptake in both
types. Especially, the CO2 uptake of the Mississippi type is more
affected by the chilled treatment than the Quebec type.

The scatter plot illustrates the positive association of ambient CO2
concentration and the CO2 uptake by the plants.

\section{Task 8:}\label{task-8}

\subsubsection{install packages and load
data}\label{install-packages-and-load-data}

\begin{Shaded}
\begin{Highlighting}[]
\CommentTok{\# install packages}
\CommentTok{\# install.packages(\textquotesingle{}devtools\textquotesingle{})}
\CommentTok{\# devtools::install\_github("hirscheylab/tidybiology")}

\CommentTok{\# load data}
\FunctionTok{library}\NormalTok{(tidyverse)}
\end{Highlighting}
\end{Shaded}

\begin{verbatim}
## Warning: package 'tibble' was built under R version 4.4.3
\end{verbatim}

\begin{verbatim}
## Warning: package 'tidyr' was built under R version 4.4.3
\end{verbatim}

\begin{verbatim}
## Warning: package 'readr' was built under R version 4.4.3
\end{verbatim}

\begin{verbatim}
## Warning: package 'lubridate' was built under R version 4.4.3
\end{verbatim}

\begin{verbatim}
## -- Attaching core tidyverse packages ------------------------ tidyverse 2.0.0 --
## v dplyr     1.2.1     v readr     2.2.0
## v forcats   1.0.1     v stringr   1.6.0
## v lubridate 1.9.5     v tibble    3.3.1
## v purrr     1.2.2     v tidyr     1.3.2
## -- Conflicts ------------------------------------------ tidyverse_conflicts() --
## x dplyr::filter() masks stats::filter()
## x dplyr::lag()    masks stats::lag()
## i Use the conflicted package (<http://conflicted.r-lib.org/>) to force all conflicts to become errors
\end{verbatim}

\begin{Shaded}
\begin{Highlighting}[]
\FunctionTok{library}\NormalTok{(tidybiology)}
\FunctionTok{data}\NormalTok{(}\StringTok{\textquotesingle{}chromosome\textquotesingle{}}\NormalTok{)}
\FunctionTok{data}\NormalTok{(}\StringTok{\textquotesingle{}proteins\textquotesingle{}}\NormalTok{)}
\end{Highlighting}
\end{Shaded}

\subsubsection{a. extract summary statistics in `chromosome'
data}\label{a.-extract-summary-statistics-in-chromosome-data}

\begin{Shaded}
\begin{Highlighting}[]
\CommentTok{\# I\textquotesingle{}m not using tidyverse functions here, cause it\textquotesingle{}s easier for me to find a concise solution in base R.}
\NormalTok{cols\_chr }\OtherTok{=} \FunctionTok{c}\NormalTok{(}\StringTok{\textquotesingle{}variations\textquotesingle{}}\NormalTok{, }\StringTok{\textquotesingle{}protein\_codinggenes\textquotesingle{}}\NormalTok{, }\StringTok{\textquotesingle{}mi\_rna\textquotesingle{}}\NormalTok{)}
\NormalTok{summary\_chr }\OtherTok{=} \FunctionTok{apply}\NormalTok{(chromosome[, cols\_chr], }\DecValTok{2}\NormalTok{, }\ControlFlowTok{function}\NormalTok{(x) \{}
  \FunctionTok{c}\NormalTok{(}\AttributeTok{mean=}\FunctionTok{mean}\NormalTok{(x, }\AttributeTok{na.rm=}\NormalTok{T),}
    \AttributeTok{median=}\FunctionTok{median}\NormalTok{(x, }\AttributeTok{na.rm=}\NormalTok{T),}
    \AttributeTok{maximum=}\FunctionTok{max}\NormalTok{(x, }\AttributeTok{na.rm=}\NormalTok{T))}
\NormalTok{\})}
\NormalTok{summary\_chr}
\end{Highlighting}
\end{Shaded}

\begin{verbatim}
##         variations protein_codinggenes    mi_rna
## mean       6484572            849.9583  73.16667
## median     6172346            836.0000  75.00000
## maximum   12945965           2058.0000 134.00000
\end{verbatim}

\subsubsection{b. visualization of chromosome size
distribution}\label{b.-visualization-of-chromosome-size-distribution}

\begin{Shaded}
\begin{Highlighting}[]
\FunctionTok{library}\NormalTok{(ggplot2)}
\FunctionTok{library}\NormalTok{(patchwork)}
\FunctionTok{ggplot}\NormalTok{(}\AttributeTok{data=}\NormalTok{chromosome, }\FunctionTok{aes}\NormalTok{(}\AttributeTok{x=}\NormalTok{id, }\AttributeTok{y=}\NormalTok{length\_mm)) }\SpecialCharTok{+} \FunctionTok{geom\_col}\NormalTok{()}
\end{Highlighting}
\end{Shaded}

\pandocbounded{\includegraphics[keepaspectratio]{assignment_1_files/figure-latex/unnamed-chunk-17-1.pdf}}

\subsubsection{c.~visualization of
relationships}\label{c.-visualization-of-relationships}

\begin{Shaded}
\begin{Highlighting}[]
\NormalTok{p1 }\OtherTok{=} \FunctionTok{ggplot}\NormalTok{(chromosome, }\FunctionTok{aes}\NormalTok{(}\AttributeTok{x=}\NormalTok{length\_mm, }\AttributeTok{y=}\NormalTok{protein\_codinggenes)) }\SpecialCharTok{+}
  \FunctionTok{geom\_point}\NormalTok{()}
\NormalTok{p2 }\OtherTok{=} \FunctionTok{ggplot}\NormalTok{(chromosome, }\FunctionTok{aes}\NormalTok{(}\AttributeTok{x=}\NormalTok{length\_mm, }\AttributeTok{y=}\NormalTok{mi\_rna)) }\SpecialCharTok{+}
  \FunctionTok{geom\_point}\NormalTok{()}

\NormalTok{p1 }\SpecialCharTok{+}\NormalTok{ p2}
\end{Highlighting}
\end{Shaded}

\pandocbounded{\includegraphics[keepaspectratio]{assignment_1_files/figure-latex/unnamed-chunk-18-1.pdf}}

Based on the plots, the chromosome size showed positive correlations
with both the number of protein coding genes, and that of miRNAs.

\subsubsection{d.~summarize the `protein'
data}\label{d.-summarize-the-protein-data}

\begin{Shaded}
\begin{Highlighting}[]
\NormalTok{cols\_protein }\OtherTok{=} \FunctionTok{c}\NormalTok{(}\StringTok{\textquotesingle{}length\textquotesingle{}}\NormalTok{, }\StringTok{\textquotesingle{}mass\textquotesingle{}}\NormalTok{)}
\NormalTok{summary\_protein }\OtherTok{=} \FunctionTok{apply}\NormalTok{(proteins[, cols\_protein], }\DecValTok{2}\NormalTok{, }\ControlFlowTok{function}\NormalTok{(x) \{}
  \FunctionTok{c}\NormalTok{(}\AttributeTok{mean=}\FunctionTok{mean}\NormalTok{(x, }\AttributeTok{na.rm=}\NormalTok{T),}
    \AttributeTok{median=}\FunctionTok{median}\NormalTok{(x, }\AttributeTok{na.rm=}\NormalTok{T),}
    \AttributeTok{maximum=}\FunctionTok{max}\NormalTok{(x, }\AttributeTok{na.rm=}\NormalTok{T))}
\NormalTok{\})}
\NormalTok{summary\_protein}
\end{Highlighting}
\end{Shaded}

\begin{verbatim}
##             length       mass
## mean      557.1603   62061.38
## median    414.0000   46140.50
## maximum 34350.0000 3816030.00
\end{verbatim}

\begin{Shaded}
\begin{Highlighting}[]
\FunctionTok{ggplot}\NormalTok{(proteins, }\FunctionTok{aes}\NormalTok{(}\AttributeTok{x=}\NormalTok{length, }\AttributeTok{y=}\NormalTok{mass)) }\SpecialCharTok{+}
  \FunctionTok{geom\_point}\NormalTok{(}\AttributeTok{alpha=}\FloatTok{0.5}\NormalTok{) }\SpecialCharTok{+}
  \FunctionTok{geom\_smooth}\NormalTok{(}\AttributeTok{method=}\StringTok{\textquotesingle{}lm\textquotesingle{}}\NormalTok{, }\AttributeTok{color=}\StringTok{"salmon"}\NormalTok{, }\AttributeTok{se=}\NormalTok{T) }\SpecialCharTok{+}
  \FunctionTok{theme\_classic}\NormalTok{()}
\end{Highlighting}
\end{Shaded}

\begin{verbatim}
## `geom_smooth()` using formula = 'y ~ x'
\end{verbatim}

\pandocbounded{\includegraphics[keepaspectratio]{assignment_1_files/figure-latex/unnamed-chunk-20-1.pdf}}

\end{document}
